\section{INTRODUCTION}

Despite rapid progress in autonomous driving, achieving reliable decision-making in long-tail and safety-critical scenarios remains a fundamental challenge \cite{kuznietsov2024explainable}. Although recent systems demonstrate strong performance within restricted operational domains, failures still occur under distribution shift, including rare road user behaviors and extreme environmental conditions \cite{xu2025wod}.
Reinforcement learning (RL) provides a principled framework for sequential decision-making through interaction-driven optimization \cite{zhao2024survey}. However, its application to autonomous driving is severely constrained by data efficiency. Learning robust driving policies typically requires massive environment interactions, yet large-scale real-world data collection is costly, time-consuming, and potentially unsafe. While high-fidelity simulators reduce safety risks, they remain computationally expensive, often requiring millions of interaction steps for policy convergence. This interaction bottleneck fundamentally limits the scalability of pure model-free RL in complex driving environments.

%To improve sample efficiency in reinforcement learning, world models (WM) learn latent representations of environment dynamics and enable policy optimization through imagination rollouts rather than repeated real-environment interaction. Building on this idea, Fig.~\ref{fig:overview} illustrates the proposed task-relevant world-model framework. Visual observations and vehicle kinematic states are integrated to form a kinematics-aware latent representation of the driving scene, while structured spatial supervision regularizes latent dynamics toward driving-relevant semantics. Policy learning is then carried out through imagination rollouts in latent space, enabling data-efficient learning with reduced dependence on real-environment interaction.

To address the interaction bottleneck, world models (WM) learn compact latent representations of environment dynamics, enabling policy optimization through imagination rollouts instead of repeated real-environment interaction. By internalizing transition dynamics in latent space, WMs significantly improve sample efficiency and support long-horizon reasoning.
Building on this paradigm, we propose a task-relevant world-model framework tailored for autonomous driving. Unlike purely generative latent modeling, our approach explicitly aligns latent dynamics with spatial and kinematic structures critical for driving tasks. Visual observations are fused with vehicle kinematic states to construct a kinematics-grounded latent representation of the driving scene, while structured spatial supervision regularizes latent transitions toward geometrically consistent and driving-relevant semantics. Policy learning is then performed via imagination rollouts in the structured latent space, enabling data-efficient and dynamically coherent decision-making with reduced reliance on real-environment interaction.

%To address the interaction bottleneck, world models (WM) learn compact latent representations of environment dynamics, enabling policy optimization through imagination rollouts instead of repeated real-environment interaction. By internalizing transition dynamics in latent space, WMs significantly improve sample efficiency and support long-horizon reasoning.
%Building on this paradigm, Fig.~\ref{fig:overview} presents the proposed task-relevant world-model framework for autonomous driving. Unlike purely generative latent modeling, our approach explicitly aligns latent dynamics with control-relevant spatial and kinematic structures. Visual observations are fused with vehicle kinematic states to construct a kinematics-grounded latent representation of the driving scene, while structured spatial supervision regularizes latent transitions toward geometrically consistent and driving-critical semantics. Policy learning is then performed via imagination rollouts in the structured latent space, enabling data-efficient and dynamically coherent decision-making with reduced reliance on real-environment interaction.


%\begin{figure}[t]
%    \centering  
%    \includegraphics[width=1.0\columnwidth]{figures/overview.png}  
%    \caption{Overview of the proposed task-relevant world-model framework.} 
%    \label{fig:overview} 
%\end{figure}

Existing approaches to autonomous driving decision-making broadly fall into several paradigms. Model-based control has long served as a foundation, where optimization-based methods such as model predictive control (MPC) explicitly exploit system dynamics and constraints for trajectory planning and tracking \cite{falcone2007predictive}. In contrast, model-free RL learns policies directly from interaction data, with algorithms such as Proximal Policy Optimization (PPO) \cite{schulman2017proximal} and Soft Actor-Critic (SAC) \cite{haarnoja2018soft} enabling stable policy optimization.
Bridging control and learning, planning with learned dynamics integrates neural networks with search or model-based rollouts \cite{silver2016mastering,ha2018recurrent}. To further improve data efficiency, WM learn compact latent representations of environment dynamics, supporting planning or imagination-based policy optimization. Representative approaches include PlaNet and Dreamer \cite{hafner2019learning,hafner2019dream,hafner2025mastering}, which perform behavior learning via latent imagination, as well as large-scale generative world models for interactive environment synthesis \cite{bruce2024genie}.
Building on these advances, world-model-based methods have been applied to autonomous driving. Prior works enhance robustness through semantic masking or representation alignment between raw observations and privileged features \cite{gao2024enhance,yang2025raw2drive}. However, existing approaches primarily focus on representation filtering or generative modeling, without explicitly constraining latent dynamics to maintain geometrically consistent and physically grounded structures essential for closed-loop vehicle control.

%Existing approaches to autonomous driving decision-making can be broadly categorized into several paradigms. Model-based control has long been a cornerstone for autonomous driving, where optimization-based controllers such as model predictive control (MPC) explicitly leverage system dynamics and constraints for motion planning and tracking; Falcone et al. \cite{falcone2007predictive} introduced predictive active steering control for vehicle stability and trajectory tracking. In contrast, model-free RL learns policies directly from interaction data; Schulman et al. \cite{schulman2017proximal} proposed Proximal Policy Optimization (PPO) as a stable on-policy baseline, while Haarnoja et al. \cite{haarnoja2018soft} introduced Soft Actor-Critic (SAC) to improve stability and performance via maximum-entropy objectives. Planning with learned models further bridged control and learning paradigms; Silver et al. \cite{silver2016mastering} demonstrated the effectiveness of combining neural networks with tree search, and Ha et al. \cite{ha2018recurrent} extended this idea through learned model-based planning. To bridge the gap between data efficiency and scalable learning, world models learn compact latent dynamics to support planning or imagination-based policy optimization: Hafner et al.  proposed PlaNet\cite{hafner2019learning} for latent dynamics planning from pixels, and further introduced Dreamer\cite{hafner2019dream} and DreamerV3\cite{hafner2025mastering} to learn behaviors through latent imagination. Recently, large-scale generative world models have pushed this paradigm toward foundation-level interactive environment generation; Bruce et al. \cite{bruce2024genie} presented Genie for generative interactive environments from visual data. Building on these developments, world-model-based approaches have been explored for autonomous driving, where Gao et al. \cite{gao2024enhance} employed semantic masking to suppress task-irrelevant visual redundancy and improve robustness, and Yang et al. \cite{yang2025raw2drive} proposed a dual-stream model based reinforcement learning approach that aligns raw sensor observations with privileged representations. However, these methods do not regularize latent dynamics toward geometrically consistent and physically grounded structures required for reliable vehicle control.

%Motivated by these observations, this paper proposes a task-relevant world-model framework that explicitly enhances spatial intelligence in autonomous driving. We introduce two complementary mechanisms. First, low-dimensional vehicle kinematic states obtained from onboard sensing are incorporated into observation encoding to ground latent dynamics in physically meaningful motion information. Second, auxiliary prediction heads are introduced to estimate lane-relative geometry and neighboring vehicle states, and their gradients are propagated back into the world model to shape latent representations toward task-critical spatial structures.
%Motivated by these observations, we propose a task-relevant world-model framework that explicitly enforces spatial and kinematic consistency in latent dynamics for autonomous driving. Rather than treating latent representations as purely generative abstractions, our approach aligns them with control-relevant structure.
%Specifically, we introduce two complementary mechanisms. First, low-dimensional vehicle kinematic states obtained from onboard sensing are incorporated into the observation encoder to ground latent transitions in physically meaningful motion dynamics. Second, structured spatial supervision is imposed through auxiliary prediction heads that estimate lane-relative geometry and neighboring vehicle states. The resulting gradients regularize latent transitions to preserve geometrically coherent and interaction-aware representations, thereby shaping the world model toward task-critical spatial structure.

Motivated by these observations, we build upon the Recurrent State-Space Model (RSSM) \cite{hafner2025mastering} and propose a task-relevant world-model framework that explicitly enforces spatial and kinematic consistency in latent dynamics for autonomous driving. Rather than treating latent representations as purely generative abstractions, our approach aligns the latent transition with driving-relevant spatial structure.
Specifically, we introduce two complementary mechanisms. First, low-dimensional vehicle kinematic states obtained from onboard sensing are incorporated into the observation encoder to ground latent transitions in physically meaningful motion dynamics. Second, structured spatial supervision is imposed through auxiliary prediction heads that estimate lane-relative geometry and neighboring vehicle states. The resulting gradients regularize the RSSM latent dynamics to preserve geometrically coherent and interaction-aware representations, thereby shaping the world model toward task-critical spatial structure.

%The contributions of this paper are summarized as follows:
%\begin{enumerate}
%    \item We propose a kinematics-aware world-model framework for autonomous driving that aligns latent dynamics with decision-critical spatial and kinematic structures.

%    \item We instantiate this within the RSSM architecture through kinematics-aware encoding that grounds perception in physical dynamics, and structured supervision that distills driving semantics into latent states via lane and vehicle prediction heads.

%    \item We demonstrate substantial improvements in data efficiency and driving performance through comprehensive experiments, alongside enhanced prediction and imagination quality of the learned world model.
%\end{enumerate}

%The contributions of this paper are summarized as follows:
%\begin{enumerate}
%    \item We propose a kinematics-grounded world-model framework for autonomous driving that explicitly aligns latent dynamics with decision-critical spatial and motion structure.

%    \item We introduce structured spatial regularization into latent dynamics learning by integrating vehicle kinematic information and geometry-aware supervision, thereby shaping the world model toward physically meaningful and interaction-aware representations.

%    \item We empirically demonstrate significant improvements in data efficiency and driving performance, together with enhanced prediction accuracy and imagination fidelity of the learned latent dynamics.
%\end{enumerate}

The contributions of this paper are summarized as follows:
\begin{enumerate}
    \item We propose a kinematics-grounded world-model framework for autonomous driving that explicitly aligns latent dynamics with decision-critical spatial and motion structure.

    \item We introduce kinematic grounding and geometry-aware spatial regularization into RSSM-based latent transition learning, guiding the latent dynamics toward physically meaningful and interaction-aware representations.

    \item We empirically demonstrate significant improvements in data efficiency and driving performance, together with enhanced prediction accuracy and imagination fidelity of the learned latent dynamics.
\end{enumerate}

The remainder of this paper is organized as follows. Section II explains our modeling problems and the preliminary knowledge. Section III details our proposed framework. Section IV introduces the experimental setup and analyzes the results. Section V concludes the paper and discusses future directions.